In a homogeneous and isotropic universe, the matter (baryonic matter + dark
matter) density parameter
$\Omega_m = \dfrac{\rho_m}{\rho_c} = 32\,\%$, where $\rho_m$ is the matter
density and $\rho_c$ is the critical density of the Universe.

\begin{description}
\item[(1)] Calculate the average matter density in our local neighbourhood.
\item[(2)] Calculate the escape velocity of a galaxy 100\,Mpc away from us.
  Assume that the recession velocity of galaxies in Hubble's law equals the
  corresponding escape velocity at that distance, for the critical density of
  the Universe that we observe.
\item[(3)] The particular galaxy is orbiting around the centre of our cluster
  of galaxies on a circular orbit. What is the angular velocity of this galaxy
  on the sky?
\item[(4)] Will we ever discriminate two such galaxies that are initially at
  the same line of sight, if they are both moving on circular orbits but at
  different radii (answer ``Yes'' or ``No'')? [Assume that the Earth is
  located at the centre of our local cluster.]
\end{description}
